\section{Initial canonical recoding under a zero-coordinate embedding}

\begin{definition}[The full ambient zero-coordinate record]
\label{hir:disjoint-zero-coordinate-initial-record}
\uses{hir:canonical-centres}
\lean{TranscendentalNakamaye.Pending.HironakaDisjointZeroCoordinateInitialRecoding.baseMarked,TranscendentalNakamaye.Pending.HironakaDisjointZeroCoordinateInitialRecoding.initialEmbedding}
\leanok
Let $n>0$, let $W\subset\mathbf C^n$ be analytic, and let $I=\mathcal I_W$ be
its whole vanishing ideal. Write $M=(\mathbf C^n,\mathbf C^n,\varnothing,I,1)$.
Let $s(x)=(x,0)$ and $p(x,t)=x$. The stabilized record is the original
$J=(\mathbf C^{n+1},\mathbf C^{n+1},\varnothing,\mathcal I_{s(W)},1)$;
in particular its smooth carrier is the full ambient manifold.
For each recursive value type define
\[
 \iota_q(v)=(1,\operatorname{inr}((0,v)))\in\operatorname{ValTy}_{q+1},
\]
where both displayed finite values retain their finite inclusions into the
top extensions.
\end{definition}

\begin{lemma}[Stabilized order and the nonspecial initial regime]
\label{hir:disjoint-zero-coordinate-initial-regime}
\uses{hir:disjoint-zero-coordinate-initial-record,hir:actual-zero-coordinate-genuine-contact-seed,hir:actual-zero-coordinate-whole-contact-lower-object}
\lean{TranscendentalNakamaye.Pending.HironakaDisjointZeroCoordinateInitialRecoding.stabilized_isInvDomain,TranscendentalNakamaye.Pending.HironakaDisjointZeroCoordinateInitialRecoding.stabilized_cosupport,TranscendentalNakamaye.Pending.HironakaDisjointZeroCoordinateInitialRecoding.stabilized_initial_orderMask,TranscendentalNakamaye.Pending.HironakaDisjointZeroCoordinateInitialRecoding.stabilized_order_one,TranscendentalNakamaye.Pending.HironakaDisjointZeroCoordinateInitialRecoding.stabilized_not_special}
\leanok
Suppose $M$ belongs to the invariant domain of dimension at most $q$.
Then $J$ belongs to the domain of dimension at most $q+1$. Its cosupport is
$s(\operatorname{cosupp}M)$, its nonmonomial order is one on that cosupport,
and its initial level is zero with no special points. Its initial order-one
mask is the whole stabilized ambient manifold.
\end{lemma}
\begin{proof}
The original graph-ideal calculation gives
$\mathcal I_{s(W)}=p^*I+(t)$ and $s^*\mathcal I_{s(W)}=I$.
The coordinate $t$ belongs to this ideal and has normal derivative one.
Thus the ideal is coherent, has order at most one everywhere, and is
nonzero at every ambient point. A cosupport point has $t=0$; its remaining
vanishing conditions are exactly those of $I$ at its projection. Conversely,
these conditions imply vanishing of every generator of $p^*I+(t)$, hence
of their entire ideal. Because the boundary is empty, the nonmonomial
order is the ordinary order, and it is exactly one on the cosupport.
The dimension condition on $M$ gives $n\le q$; the full stabilized carrier
therefore has dimension $n+1\le q+1$. Empty boundaries are simple normal
crossings and transverse to each full carrier, proving the domain claim.
There are no original boundary members, so the initial level is zero.
If a point belonged to the codimension-one part of the stabilized
cosupport, the maximal-contact equation $t$ would identify that cosupport
locally with the entire plane $t=0$. Every original germ in $I$ at the
projected point would consequently vanish on a neighborhood in
$\mathbf C^n$: pull it back by $p$, use its vanishing on the stabilized
cosupport, then restrict by $s$. This contradicts the nonzero-along-carrier
condition in the original domain. The positive-level special case is
impossible because the level is zero. Thus there are no special points.
\end{proof}

\begin{lemma}[The actual initial contact value]
\label{hir:disjoint-zero-coordinate-initial-contact-value}
\uses{hir:disjoint-zero-coordinate-initial-regime,hir:actual-zero-coordinate-stepone-object,hir:actual-zero-coordinate-whole-contact-lower-object,hir:actual-closed-section-all-index-initial-values,mc:contact-lower-value}
\lean{TranscendentalNakamaye.Pending.HironakaDisjointZeroCoordinateInitialRecoding.stabilized_stepOne_initial_value}
\leanok
At every $x\in\operatorname{cosupp}M$, Step One for $J$, using the actual
inductive canonical family of depth $q$ in ambient dimension $n+1$, has
initial value $(0,\operatorname{inv}_q(M,x))$.
\end{lemma}
\begin{proof}
The marking-one homogenized and coefficient records of the stabilized
ideal are the original stabilized record. The actual normal coordinate
$t$ is a global maximal-contact direction. Its coefficient record on
$H=\{t=0\}$ is literally
$(\mathbf C^{n+1},H,\varnothing,\mathcal I_{s(W)},1)$; the whole ideal
restricts by $s$ to $I$, with no choice of generators imposed on the invariant.
This record is in the depth-$q$ domain: its carrier has dimension $n\le q$,
and projection and section identify its restricted ideal with the original
nonzero ideal $I$. Restrict both ambient manifolds to their whole open sets
and apply the existing initial closed-section value theorem to this literal
record. The identity $s^*\mathcal I_{s(W)}=I$ then gives its depth-$q$ value
exactly as $\operatorname{inv}_q(M,x)$.
The actual independently chosen Step One direction can be replaced by $t$
by the maximal-contact reference-independence theorem. Its required lower
domain is the preceding coefficient domain; the chosen lower domain and
cosupport follow from the nonspecial initial regime. Every relevant tower
and restricted tail is empty, so canonicity is automatic and both run
clocks are zero. The whole-open restriction property of the invariant
then gives precisely the stated finite lower value. The outer Step One
level is zero by the preceding lemma.
\end{proof}

\begin{theorem}[Full ambient initial depth-shift law]
\label{hir:disjoint-zero-coordinate-initial-value}
\uses{hir:disjoint-zero-coordinate-initial-contact-value,hir:disjoint-zero-coordinate-initial-record}
\lean{TranscendentalNakamaye.Pending.HironakaDisjointZeroCoordinateInitialRecoding.stabilized_initial_value}
\leanok
For the actual original canonical family and every
$x\in\operatorname{cosupp}M$,
\[
 \operatorname{inv}_{q+1}(J,s(x))
 =\iota_q\bigl(\operatorname{inv}_q(M,x)\bigr).
\]
\end{theorem}
\begin{proof}
The first component of the general invariant is the nonmonomial order,
which is one by the stabilized-order lemma. Its initial order mask is the
whole ambient manifold. With empty boundary and mark one, the initial
companion on this mask is restricted-equivalent to the stabilized marked
ideal itself: the monomial factor is the unit ideal and the nonmonomial
factor is its original ideal restricted to the carrier. The actual
maximal-order Step One invariant is unchanged under this representative
change and under the whole-open identification. Its value is the finite
pair computed in the preceding lemma. The definition of the general
invariant therefore gives exactly the finite right-branch expression
$\iota_q(\operatorname{inv}_q(M,x))$.
\end{proof}

\begin{theorem}[The first full ambient maximum locus]
\label{hir:disjoint-zero-coordinate-initial-centre}
\uses{hir:disjoint-zero-coordinate-initial-value}
\lean{TranscendentalNakamaye.Pending.HironakaDisjointZeroCoordinateInitialRecoding.stabilized_initial_maximum_locus,TranscendentalNakamaye.Pending.HironakaDisjointZeroCoordinateInitialRecoding.stabilized_canonical_initial_maximum_locus}
\leanok
The initial maximum locus of the full ambient depth-$(q+1)$ invariant of
$J$ is the zero-section image of the initial maximum locus of the original
depth-$q$ invariant of $M$. When $q=n$, this is the first-centre comparison
for the actual canonical invariants in ambient dimensions $n+1$ and $n$.
\end{theorem}
\begin{proof}
The map $\iota_q$ preserves and reflects order: every added coordinate and
branch is constant, and the finite inclusions in top extensions preserve
and reflect the remaining order. The stabilized cosupport is exactly the
zero-section image of the original one. At each of its points substitute
the depth-shift law into the defining comparison with every other
cosupport point. Reflection of order gives exactly the old maximum
condition at the projected point; preservation of order proves the
converse. This argument also applies when the cosupport is empty: both
maximum loci are then empty by their definitions. For $q=n>0$, the
canonical invariant is definitionally the corresponding inductive family;
the same identification in dimension $n+1$ gives the claimed canonical
first-centre formula.
\end{proof}

\begin{lemma}[Empty cosupport is preserved]
\label{hir:disjoint-zero-coordinate-initial-empty}
\uses{hir:disjoint-zero-coordinate-initial-regime}
\lean{TranscendentalNakamaye.Pending.HironakaDisjointZeroCoordinateInitialRecoding.stabilized_cosupport_eq_empty_iff}
\leanok
The stabilized cosupport is empty if and only if the original cosupport is
empty.
\end{lemma}
\begin{proof}
The stabilized cosupport is the zero-section image of the original one.
A point of the latter gives its section point in the former. Conversely,
the projection of any point of the former belongs to the latter. Thus one
is empty exactly when the other is empty.
\end{proof}

\begin{lemma}[The whole base space is excluded by the natural domain]
\label{hir:disjoint-zero-coordinate-initial-proper}
\uses{hir:disjoint-zero-coordinate-initial-record}
\lean{TranscendentalNakamaye.Pending.HironakaDisjointZeroCoordinateInitialRecoding.base_domain_ne_univ}
\leanok
The original invariant-domain hypothesis implies $W\ne\mathbf C^n$.
\end{lemma}
\begin{proof}
If $W=\mathbf C^n$, every germ in its vanishing ideal vanishes on a whole
neighborhood. At the origin this contradicts the required existence of a
member germ nonzero arbitrarily nearby.
\end{proof}

\begin{remark}
The theorem includes $W=\varnothing$ and every other empty-cosupport case.
It is an initial-state law. It does not assert stabilization of finite
histories or intrinsic descent of a complete resolution algorithm.
\end{remark}
